\subsection{Deferred intervention measures}\label{app:deferredmeasures}

The study excludes intervention-rate outcomes but retains help and
administrative records for eligibility and cost. These definitions are
proposed future measures, not collected results.
HG-DAgger and ThriftyDAgger supply cross-domain measurement precedents,
not coding-agent effectiveness evidence
~\cite{kelly2019hgdagger,hoque2022thriftydagger}.

For run $r$, $n_r$ is attempted feature cycles, including assigned
post-steps and failed cycles, and $h_{rj}$ qualifying corrective-help
events. Observed human intervention rate would be
\begin{equation}
\mathrm{HIR}_r=\frac{1}{n_r}\sum_{j=1}^{n_r}
\mathbf{1}\{h_{rj}>0\}.
\label{eq:hir}
\end{equation}
Event counts and active human minutes are separate; repeated help in one
cycle counts once for incidence. With $d_{rj}=1$ for an explicit help
request or a block needing corrective human action,
\begin{equation}
\mathrm{IDR}_r=\frac{1}{n_r}\sum_{j=1}^{n_r}d_{rj}.
\label{eq:intervention-demand}
\end{equation}
Requests and supplied help differ. A timeout alone is not a help request;
zero attempts gives undefined rates. A future annotation protocol needs
event/cycle/phase locators, coverage checks and adjudication; absent
records cannot be zero. Routine launches, planned approvals and safety
stops differ from corrective rescue. Low intervention must not reward
silent failure or discourage necessary human judgement.

\subsection{Historical freeze provenance}\label{app:freezeprovenance}

The reported case-tool snapshot is HVE Core \texttt{plugins-v3.3.106},
snapshot commit
\nolinkurl{2df5013e29ebfd2ce61bcd0cf021d8c8186482eb},
generated from source
\nolinkurl{e9af018fe5ffbfbc72321a8bca21fb987359af08}.
These are reported provenance identifiers, not verification of installation
date or every session's runtime.

The historical v4 execution source was
\nolinkurl{e8707ecc8965d3f81f5142a2603042bfca1aab15}, with
manifest-only commit
\nolinkurl{03c50856838b539780d57216af2a4cf6744187d7}.
Its attached \nolinkurl{osf-v4-2026-10-31.tar.gz} checksum was
\nolinkurl{1a42fac40a4f1aa82df59c8c3f7f8d6dce37d22502c5952c398362ac63d6bd65}.
The runtime image was
\nolinkurl{sha256:090220b0b172ec205f1ea473d7c3a223eec5d7d619eb21acdcf15c9289287fa3}.
These identify historical inputs, not a freeze or approval of the
five-arm amendment. The original registration at
\url{https://osf.io/ywsrv/} is owner-reported approved after submission
on 28 October; the inspected v4 prospective amendment flag was false
and its completion date null. New approval and actual prospective
recording cannot be inferred or backdated.
